% Pitch del proyecto final. 5 minutos, 11 diapositivas, unos 30 segundos cada una.
% Se puede usar otra herramienta si conserva este orden y se entrega como pitch.pdf.
% Las cifras deben coincidir con 02_finanzas/modelo_financiero.xlsx y con el artículo.
% Regla de diseño: una idea por diapositiva, poco texto, letra grande, una imagen.
% Compilar con: pdflatex pitch.tex
\documentclass[aspectratio=169,11pt]{beamer}
\usepackage[utf8]{inputenc}
\usepackage[T1]{fontenc}
\usepackage[spanish]{babel}
\decimalpoint
\usetheme{metropolis}
\usepackage{booktabs,xcolor}
\newcommand{\pendiente}[1]{\textcolor{gray}{\emph{#1}}}

\title{\pendiente{Nombre de la empresa}}
\subtitle{\pendiente{La empresa en una oración: qué hace y para quién}}
\author{\pendiente{Integrantes}}
\date{}

\begin{document}

% 1. Propósito
\begin{frame}[plain]
  \vspace*{-15.7pt}  % compensa el desfase vertical de la portada de metropolis
  \titlepage
\end{frame}

% 2. Problema: una persona concreta y su dolor, con una cifra.
\begin{frame}{\pendiente{El dolor del cliente, en una frase}}
  \pendiente{Quién sufre el problema, cada cuánto y cuánto le cuesta. Una cita de entrevista.}
\end{frame}

% 3. Tamaño de mercado: TAM, SAM y SOM de la hoja Resultados (filas 22 a 25).
\begin{frame}{\pendiente{Tamaño del mercado}}
  \centering
  \begin{tabular}{lrr}
    \toprule
    & Clientes & MXN al año \\
    \midrule
    TAM & \pendiente{0} & \pendiente{0} \\
    SAM & \pendiente{0} & \pendiente{0} \\
    SOM & \pendiente{0} & \pendiente{0} \\
    \bottomrule
  \end{tabular}

  \vspace{1em}
  \pendiente{Fuente del número de clientes potenciales.}
\end{frame}

% 4. Solución: qué ofrece la empresa, sin tecnicismos.
\begin{frame}{\pendiente{La solución}}
  \pendiente{Qué recibe el cliente y qué deja de hacer.}
\end{frame}

% 5. Producto: demostración, arquitectura y el papel de la IA, con el resultado del artículo.
\begin{frame}{\pendiente{El producto}}
  \pendiente{Captura o video del MVP. Dónde entra la IA. Resultado principal contra la línea base.}
\end{frame}

% 6. Validación: entrevistas, hipótesis validadas e invalidadas, métricas de la landing.
\begin{frame}{\pendiente{Lo que aprendimos de los clientes}}
  \pendiente{Número de entrevistas, hipótesis validadas e invalidadas, visitas y registros.}
\end{frame}

% 7. Competencia: matriz con dos atributos que importan al cliente.
\begin{frame}{\pendiente{Competencia}}
  \pendiente{Matriz de dos ejes o tabla con competidores y la ventaja de la empresa.}
\end{frame}

% 8. Modelo de negocio: cómo se gana dinero.
\begin{frame}{\pendiente{Modelo de negocio}}
  \pendiente{Qué se cobra y cuánto. Canal de adquisición. CAC, LTV y LTV/CAC.}
\end{frame}

% 9. Finanzas: corrida a 24 meses, capital necesario y retorno.
\begin{frame}{\pendiente{Finanzas a 24 meses}}
  \pendiente{Gráfica del flujo acumulado. Capital necesario, mes de equilibrio, VPN. Supuestos principales.}
\end{frame}

% 10. Equipo: nombres, fotos y roles.
\begin{frame}{\pendiente{Equipo}}
  \pendiente{Foto, nombre y rol de cada integrante.}
\end{frame}

% 11. Cierre: lo que la empresa pide (capital, piloto, contactos) y la frase de la portada.
\begin{frame}{\pendiente{Lo que necesitamos}}
  \pendiente{Monto o recurso que se pide y para qué. Repetir la frase de la empresa.}
\end{frame}

\end{document}
